%%%%%%%%%%  scan idx 376  |  folio 363  %%%%%%%%%%
un homomorphisme de groupes de \(\widehat G(K)\) dans \(K[[t]]\); l'application qui à une série associe son \((k-1)\)-ième coefficient, multiplié par \((-1)^{k-1}\), est un homomorphisme additif de \(K[[t]]\) dans \(K\)~; le morphisme \(\Phi^k\) sera le composé des deux. Il est clair que \(\Phi^k\circ\lambda_t=\psi^k\) d'après la définition (7.1.1).

La seconde des formules (7.5.1) exprime donc la commutativité du diagramme écrit plus haut. Si on suppose que \(K\) est un \(\lambda\)-anneau, le carré de gauche est commutatif, et il suffit de vérifier la commutativité de celui de droite. Or il s'agit d'un diagramme d'homomorphismes de foncteurs sur \(\mathcal C_0\), ayant pour source \(\widehat G\times\widehat G\)~; il suffit donc d'après 1.10 de vérifier la commutativité pour \((1+Xt,1+Yt)\)~; on a alors :
\[
  \Phi^k(1+Xt)=X^k\qquad\bigl(\text{$(k-1)$-ième coefficient de }X/1+Xt\bigr)~;\qquad
  \Phi(1+Yt)=Y^k~;
\]
\(\Phi^k\bigl((1+Xt)\circ(1+Yt)\bigr)=\Phi^k(1+XYt)=(XY)^k=\Phi^k(1+Xt)\Phi^k(1+Yt)\), ce qui donne le résultat.

Pour montrer la dernière, on remarque d'abord que le diagramme :
\begin{equation}\tag{7.5.3}
\begin{tikzcd}[column sep=large,row sep=large]
\widehat G^t(K) \arrow[r,"\lambda_u"] \arrow[d,"\Phi^{kk'}"'] &
\widehat G^u\circ\widehat G^t(K) \arrow[d,"\Phi^{k'}_u"]\\
K \arrow[r,"\Phi^k"] &
\widehat G^t(K),
\end{tikzcd}
\end{equation}
est commutatif. Comme plus haut, il suffit en effet de le voir pour \(1+Xt\) et on a :
\[
  \Phi^{k'}_u(\lambda_u(1+Xt))
  =
  \Phi^{k'}_u(1+\{1+Xt\}u)
  =
  (1+Xt)^{\circ k'}
  =
  1+X^{k'}t~;
\]
\[
  \Phi^k(1+X^{k'}t)=X^{kk'}=\Phi^{kk'}(1+Xt).
\]

Par suite, \(\psi^{kk'}=\Phi^k\circ\Phi^{k'}_u\circ\lambda_u\circ\lambda_t=\Phi^k\circ\psi^{k'}\circ\lambda_t\).

Par ailleurs, \(\psi^k\circ\psi^{k'}=\Phi^k\circ\lambda_t\circ\Phi^{k'}\circ\lambda_t=\Phi^k\circ\lambda_t\circ\psi^{k'}\).
